Udgangspunktet for denne opgave var at formidle et økologisk budskab om, at menneskeheden skal beskytte naturen, her havet. Det findes, at især havstrømme er en stor faktor i at frembringe plastik, hvorefter disse klynger sig sammen. Dette er værd at kommunikere til en spiller, men skal gøres på et niveau, der er rimeligt. Til dette formål, kan der tages kunstneriske friheder, således at man laver tegneserieagtige havstrømme på et meget lokalt niveau, hvilket kan opnås hvis hjælp af tyngdefeltbetragtninger, gode vektorberegninger og stabile numeriske metoder som semi-implicit Eulers. Alt dette kan indpakkes i form af et antropocænt omsvøb med arkade-charme, der foremår at være appellerende til Gen Z ved hjælp af anemoia og menneskelig heroisme samtidig med at budskabet er klart i form af antagonistiske monstre og klaustrofobisk skraldeoverdynge.
